% Part: model-theory
% Chapter: models-of-arithmetic
% Section: introduction

\documentclass[../../../include/open-logic-section]{subfiles}

\begin{document}

\olfileid{mod}{mar}{int}
\section{سريزه}

د حساب \emph{معياري مدل} هغه !!{structure}~$\Struct{N}$ دے چې
$\Domain{N} = \Nat$ وي او $\Obj{0}$، $\prime$، $+$، $\times$ او $<$
هغسې تفسير شوي وي لکه تمه چې ترې کېږي۔ يعنې $\Obj{0}$، $0$ دے،
$\prime$ د جانشين تابع ده، $+$ د جمع په توګه او $\times$ د ضرب په
توګه د~$\Nat$ د شمېرو تفسير لري۔ په ځانګړي ډول:
\begin{align*}
  \Assign{\Obj{0}}{N} & = 0\\
  \Assign{\prime}{N}(n) & = n + 1\\
  \Assign{+}{N}(n, m) & = n + m\\
  \Assign{\times}{N}(n, m) & = nm
\end{align*}
البته د~$\Lang{L_A}$ داسې جوړښتونه هم شته چې تعريف ساحې ئې له~$\Nat$
څخه بېلې دي۔ د بېلګې په توګه، $\Struct{M}$ داسې ونيسئ چې
$\Domain{M} = \{a\}^*$ وي (د يوازېنۍ نښې~$a$ متناهي لړۍ؛ يعنې
$\emptyset$، $a$، $aa$، $aaa$، \dots)، او تفسيرونه ئې داسې وي:
\begin{align*}
  \Assign{\Obj{0}}{M} & = \emptyset\\
  \Assign{\prime}{M}(s) & = s \concat a\\
  \Assign{+}{M}(n, m) & = a^{n + m}\\
  \Assign{\times}{M}(n, m) & = a^{nm}
\end{align*}
دا دوه جوړښتونه په دې معنا «په اصل کښې يو شان» دي چې يوازينے
توپير ئې د !!{element}s او !!{domain}s په څيزونو کښې دے، نه دا چې
د تفسير د تابعو له مخې د !!{domain}s !!{element}s له يو بل سره څنګه تړاو لري۔
موږ وايو چې دا دوه !!{structure}s \emph{آيزومورف} دي۔

د متناهي شاهد له قضيه څخه په اسانۍ پايله کېږي چې هره تيوري چې په
$\Struct{N}$ کښې رښتونې وي، داسې مدلونه هم لري چې له~$\Struct{N}$ سره
آيزومورف نۀ وي۔ دا ډول جوړښتونه \emph{نامعياري} بلل کېږي۔ په معياري
مدل کښې (يعنې په~$\Struct{N}$ او له هغې سره په ټولو آيزومورفو
!!{structure}s کښې) ټول !!{element}s د معياري عددنښو په ارزښتونو
$\num{n}$ کښې راټولېږي، يعنې:
\[
\Domain{N} = \Setabs{\Value{\num{n}}{N}}{n \in \Nat}
\]
خو په نامعياري مدلونو کښې داسې نۀ ده: که $\Struct{M}$ نامعياري وي، نو
لږ تر لږه يو $x \in \Domain{M}$ شته چې د هېڅ~$n$ دپاره
$x \neq \Value{\num{n}}{M}$ وي۔

دا نامعياري عناصر په زړه پورې دي: «نامتناهي طبيعي شمېرې» دي۔
خو د هغو شتون د ناتکميلۍ پديدې هم په يوه معنا روښانوي۔ د بېلګې په
توګه د پيانو د حساب د سازګارۍ جمله، $\OCon[\Th{PA}]$، يعنې
$\lnot \lexists[x][\OPrf[\Th{PA}](x, \gn{\lfalse})]$، راواخلئ۔
څرنګه چې $\Th{PA}$ نه $\OCon[\Th{PA}]$ ثابتوي او نه
$\lnot \OCon[\Th{PA}]$، نو هره يوه جمله په سازګار ډول له~$\Th{PA}$
سره يوځاے کېدے شي۔ څرنګه چې $\Th{PA}$ سازګاره ده،
$\Sat{N}{\OCon[\Th{PA}]}$ او په پايله کښې
$\Sat/{N}{\lnot \OCon[\Th{PA}]}$ لرو۔ نو $\Struct{N}$ د
$\Th{PA} \cup \{\lnot \OCon[\Th{PA}]\}$ مدل \emph{نۀ} دے، او ټول
مدلونه ئې بايد نامعياري وي۔ د $\Th{PA} \cup
\{\lnot \OCon[\Th{PA}]\}$ مدلونه بايد داسې يو !!{element} ولري چې د
$\lexists[x][\OPrf[\Th{PA}](x,\gn{\lfalse})]$ د رښتينولۍ شاهد وي، يعنې
له~$\Th{PA}$ څخه د تناقض د !!a{derivation} ګوډل عدد وي۔ داسې
!!a{element} معياري کېدے نۀ شي، ځکه د هر~$n$ دپاره
$\Th{PA} \Proves \lnot \OPrf[\Th{PA}](\num{n}, \gn{\lfalse})$۔
\textbf{د ژباړې د نسخې سمون (OLMOD-004):} سرچينې د وجودي جملې په
دويم ځل ليکلو کښې د ثبوت د کوډ ارګومېنټ x غورځولے و؛ دلته د ثبوت د
نسبت دوه ارګومېنټه بڼه بېرته ليکل شوې، لکه د همدې فقرې په لومړي
تعريف او په ورپسې معياري عددنښه کښې چې راغلې ده۔

\end{document}
